% mnras_template.tex 
%
% LaTeX template for creating an MNRAS paper
%
% v3.0 released 14 May 2015
% (version numbers match those of mnras.cls)
%
% Copyright (C) Royal Astronomical Society 2015
% Authors:
% Keith T. Smith (Royal Astronomical Society)

% Change log
%
% v3.0 May 2015
%    Renamed to match the new package name
%    Version number matches mnras.cls
%    A few minor tweaks to wording
% v1.0 September 2013
%    Beta testing only - never publicly released
%    First version: a simple (ish) template for creating an MNRAS paper

%%%%%%%%%%%%%%%%%%%%%%%%%%%%%%%%%%%%%%%%%%%%%%%%%%
% Basic setup. Most papers should leave these options alone.
\documentclass[fleqn,usenatbib]{mnras}

% MNRAS is set in Times font. If you don't have this installed (most LaTeX
% installations will be fine) or prefer the old Computer Modern fonts, comment
% out the following line
%\usepackage{newtxtext,newtxmath}
% Depending on your LaTeX fonts installation, you might get better results with one of these:
%\usepackage{mathptmx}
%\usepackage{txfonts}

% Use vector fonts, so it zooms properly in on-screen viewing software
% Don't change these lines unless you know what you are doing
\usepackage[T1]{fontenc}

% Allow "Thomas van Noord" and "Simon de Laguarde" and alike to be sorted by "N" and "L" etc. in the bibliography.
% Write the name in the bibliography as "\VAN{Noord}{Van}{van} Noord, Thomas"
\DeclareRobustCommand{\VAN}[3]{#2}
\let\VANthebibliography\thebibliography
\def\thebibliography{\DeclareRobustCommand{\VAN}[3]{##3}\VANthebibliography}


%%%%% AUTHORS - PLACE YOUR OWN PACKAGES HERE %%%%%

% Only include extra packages if you really need them. Common packages are:
\usepackage{graphicx}	% Including figure files
\usepackage{amsmath}	% Advanced maths commands
\usepackage{amssymb}	% Extra maths symbols

\usepackage{tabularx}
\usepackage{comment}
\usepackage{float}
\usepackage{xcolor}

%%%%%%%%%%%%%%%%%%%%%%%%%%%%%%%%%%%%%%%%%%%%%%%%%%

%%%%% AUTHORS - PLACE YOUR OWN COMMANDS HERE %%%%%

% Please keep new commands to a minimum, and use \newcommand not \def to avoid
% overwriting existing commands. Example:
%\newcommand{\pcm}{\,cm$^{-2}$}	% per cm-squared

\newcommand{\han}{\textcolor{orange}}
\newcommand{\fref}{f_{\rm ref}}
\newcommand{\ARed}{A_{\rm sn}}
\newcommand{\gammaRed}{\gamma_{\rm sn}}
\newcommand{\ARedI}{A_{{\rm sn},I}}
\newcommand{\ARedJ}{A_{{\rm sn},J}}
\newcommand{\gammaRedI}{\gamma_{{\rm sn},I}}
\newcommand{\gammaRedJ}{\gamma_{{\rm sn},J}}
\newcommand{\Pg}{P_{\rm g}}
\newcommand{\PI}{P_I}
\newcommand{\PJ}{P_J}
\newcommand{\PsnI}{P_{{\rm sn},I}}
\newcommand{\chiIJhd}{\chi_{IJ}}
\newcommand{\chiIJdip}{\chi_{IJ}^{\rm dip}}
\newcommand{\sigJitter}{\sigma_{\rm jit}}

%%%%%%%%%%%%%%%%%%%%%%%%%%%%%%%%%%%%%%%%%%%%%%%%%%

%%%%%%%%%%%%%%%%%%% TITLE PAGE %%%%%%%%%%%%%%%%%%%

% Title of the paper, and the short title which is used in the headers.
% Keep the title short and informative.
\title[Short title, max. 45 characters]{Improving pulsar timing array sensitivity for MeerTime}

% The list of authors, and the short list which is used in the headers.
% If you need two or more lines of authors, add an extra line using \newauthor
\author[K. T. Smith et al.]{
First Author,$^{1}$\thanks{E-mail: mn@ras.org.uk (KTS)}
Second Author,$^{2}$
Third Author$^{2,3}$
and Fourth Author$^{3}$
and ,,,,,
\\
% List of institutions
$^{1}$Institution one\\
$^{2}$Institution two\\
$^{3}$Institution three
}

% These dates will be filled out by the publisher
\date{Accepted XXX. Received YYY; in original form ZZZ}

% Enter the current year, for the copyright statements etc.
\pubyear{2015}

% Don't change these lines
\begin{document}
\label{firstpage}
\pagerange{\pageref{firstpage}--\pageref{lastpage}}
\maketitle

% Abstract of the paper
\begin{abstract}
Abstract here
\end{abstract}

% Select between one and six entries from the list of approved keywords.
% Don't make up new ones.
\begin{keywords}
keyword1 -- keyword2 -- keyword3
\end{keywords}

%%%%%%%%%%%%%%%%%%%%%%%%%%%%%%%%%%%%%%%%%%%%%%%%%%

%%%%%%%%%%%%%%%%% BODY OF PAPER %%%%%%%%%%%%%%%%%%

\section{Introduction}

\section{PTA average signal to noise ratio}

Here we describe how we calculate the PTA signal to noise radio (S/N) for a set of pulsars and a predicted gravitational wave background signal. 
We follow the method in~\cite{SiemensEtAl:2013}. 

The gravitational wave spectrum is given by 
\begin{equation}
\Pg (f) = \frac{A^2}{24 \pi^2} \left(\frac{f}{\fref}\right)^{2\alpha} f^{-3}\,,
\label{eqn:Pgw}
% equation 16 in SiemansEtAl:2013
\end{equation}
where $f$ is the gravitational wave frequency, $\fref$ is a reference frequency chosen as $1\,{\rm yr}^{-1}$, $A$ is the amplitude of the background at $f=\fref$, and $\alpha$ is the slope of the spectrum. 
Throughout this work we assume $a=-2/3$ consistent with a stochastic background from circular binaries [ref] and $A = 2 \times 10^{-15}$ consistent with recent hints of a possible signal [ref].

For a gravitational wave background signal and a given pulsar timing array with timespan $T$, the average S/N is given by
\begin{equation}
\left< \rho \right> = \left( 2 T \sum_{IJ} \chiIJhd^2 
                             \int_{f_L}^{f_H} {\rm d}f  
                             \frac{\Pg^2 (f)}{\PI(f)\PJ(f)}
                            \right)^{1/2} \,,
\label{eqn:snrPIPJPg}
\end{equation}
where $\sum_{IJ}$ is the sum over all pulsar pairs ($I \neq J$) and $P_I$ is the power spectrum of the $I{\rm th}$ pulsar. 
The integration limits are given by $f_L=1/T$ and $f_H=c/2$ where $c$ is the cadence of observations, which we assume to be the identical for all pulsars in the array. 
The Hellings and Downs coefficients $\chi_{IJ}$ are given by [ref]
\begin{eqnarray}
\chiIJhd &=& \frac{1}{2} 
           - \frac{1}{4} \left(\frac{1 - \cos(\zeta)}{2}\right) \nonumber\\
           & &+ \frac{3}{2} \left(\frac{1 - \cos(\zeta)}{2}\right) 
           \log \left( \frac{1 - \cos(\zeta)}{2} \right)\,,
\label{eqn:hd}
\end{eqnarray}
for pulsars separated by angle $\zeta$.
If we assume evenly sampled data and that there is only white noise in the pulsars then $\PI$ is the sum of Eqn.~\ref{eqn:Pgw} plus white noise, 
\begin{equation}
\PI (f) = \frac{A^2}{24 \pi^2} \left(\frac{f}{\fref}\right)^{2\alpha} f^{-3} 
          + 2 \sigma_I^2 \Delta t\,,
\end{equation}
where $c=1/\Delta t$ and $\sigma_I$ is the timing precision for the $I{\rm th}$ pulsar. 

Pulsars may also contain spin noise (or red noise) and jitter noise. 
The power spectral density of spin noise for the $I{\rm th}$ pulsar is given by 
\begin{equation}
% from Goncharov 2021 eqn 5
\PsnI (f) = \frac{\ARedI^2}{12 \pi^2} 
            \left( \frac{f}{\fref} \right)^{-\gammaRedI} {\rm yr}^3\,,
\end{equation}
where $\ARedI$ and $\gammaRedI$ are the amplitude and slope of the power law for the $I$th pulsar and ${\rm yr}$ is the number of seconds in a year. 
Jitter noise $\sigma_{\rm jit}$ is added in quadrature to the white noise part the the spectrum.
The resulting power spectrum including GWB, white noise, spin noise, and jitter noise is given by
\begin{eqnarray}
\PI (f) &=& \frac{A^2}{24 \pi^2} \left(\frac{f}{\fref}\right)^{2\alpha} f^{-3} \nonumber\\
          & &~+~ 2 \left( \sigma_I^2 + \sigma_{{\rm jit}, I}^2 \right) \Delta t \nonumber\\
          & &~+~ \frac{\ARedI^2}{12 \pi^2} 
          \left( \frac{f}{\fref} \right)^{-\gammaRedI} {\rm yr}^3\,,
\label{eqn:PIGWWR}
\end{eqnarray}
\begin{comment}
\begin{equation}
\left< \rho \right> = \left( 2 T \sum_{IJ} \chi_{IJ}^2
                             \int_{f_L}^{f_H} {\rm d}f  
                             \frac{\Pg^2 (f)}
                                  {[\Pg (f) + 2\sigma_I^2\Delta t + \Psn (f)]
                                   [\Pg (f) + 2\sigma_J^2\Delta t + \Psn (f)]} 
                            \right)^{1/2} \,.
\label{eqn:snrWhiteNoiseOnly}
\end{equation}
\end{comment}


\section{Data}

MeerTime is timing $82$ pulsars and currently targets a precision of $1\times 10^{-6}\,{\rm s}$ for each with $c=26\,{\rm yr^{-1}}$.
From observations to date, we have an $\sigma_I$ and the integration time required to achieve it for each pulsar and $\sigma_I$ scales with observation $T_{{\rm obs},I}$ time like 
\begin{equation}
\sigma_I = \delta_I / \sqrt{T_{{\rm obs},I}}\,,
\end{equation}
where $\delta_I$ is a constant of proportionality which depends on the pulsar. 

Limited information on the pulsar spin and jitter noise is also available as described in the following sections. 

\subsection{Spin noise}
MeerTime observations are not currently long enough to estimate spin noise in each of the pulsars. 
However spin noise measurements already exist for $13$ of the pulsars from PPTA and NANOGrav analysis. 
We take $\ARed$ and $\gammaRed$ values from ~\cite{GoncharovEtAl:2021} (sixth and seventh column in Table 1) and ~\cite{NANOGrav12p5yr:2021} (tenth and eleventh column in Table 5). 
Where a value is listed in both papers, we take the~\cite{GoncharovEtAl:2021} value. 
Table~\ref{tab:redNoise} shows the red noise parameters for $13$ of our $82$ pulsars. 

\begin{table}
\begin{center}
\begin{tabular}{llll}
\hline \hline
Pulsar     & $\log_10(\ARed)$  & $\gammaRed$ & Reference \\
\hline
J0030+0451 & $-15.72$ & $5.3$   &  \cite{NANOGrav12p5yr:2021}\\
J0437-4715 & $-14.56$ & $2.99$  &  \cite{GoncharovEtAl:2021}\\
J0613-0200 & $-14.26$ & $4.17$  &  \cite{GoncharovEtAl:2021}\\
J0711-6830 & $-13.04$ & $1.09$  &  \cite{GoncharovEtAl:2021}\\
J1024-0719 & $-14.62$ & $6.39$  &  \cite{GoncharovEtAl:2021}\\
J1600-3053 & $-14.34$ & $3.81$  &  \cite{GoncharovEtAl:2021}\\
J1643-1224 & $-12.85$ & $0.98$  &  \cite{GoncharovEtAl:2021}\\
J1713+0747 & $-15.16$ & $1.6$   &  \cite{NANOGrav12p5yr:2021}\\
J1744-1134 & $-14.16$ & $1.7$   &  \cite{NANOGrav12p5yr:2021}\\
J1747-4036 & $-13.78$ & $3.8$   &  \cite{NANOGrav12p5yr:2021}\\
J1909-3744 & $-14.74$ & $4.05$  &  \cite{GoncharovEtAl:2021}\\
J2145-0750 & $-13.86$ & $1.6$   &  \cite{NANOGrav12p5yr:2021}\\
J2317+1439 & $-16.50$ & $6.0$   &  \cite{NANOGrav12p5yr:2021}\\
\hline\hline
\end{tabular}
\end{center}
\caption{\label{tab:redNoise}
Spin noise parameters for $13$ pulsars. 
The second and third columns show the log of the spin noise amplitude and the slope. 
}
\end{table}

% pulsar reference - jitter, white nooise, red noies, position

\begin{comment}
% value from NANOGRav table is An
% An  = 0.006 \mu s yr^{1/2} 
% An -> 0.006 * 3.17E-14 yr^{3/2}
\begin{table}
\begin{tabular}{llll}
\hline\hline
Pulsar     & $\ARed$ & $\gammaRed$ & Ref.\\
\hline
J0030+0451 & $-15.72$ & $5.3$   &  \cite{NANOGrav12p5yr:2021}\\
J0437-4715 & $-14.56$ & $2.99$  &  \cite{GoncharovEtAl:2021}\\
J0613-0200 & $-14.26$ & $4.17$  &  \cite{GoncharovEtAl:2021}\\
J0711-6830 & $-13.04$ & $1.09$  &  \cite{GoncharovEtAl:2021}\\
J1024-0719 & $-14.62$ & $6.39$  &  \cite{GoncharovEtAl:2021}\\
J1600-3053 & $-14.34$ & $3.81$  &  \cite{GoncharovEtAl:2021}\\
J1643-1224 & $-12.85$ & $0.98$  &  \cite{GoncharovEtAl:2021}\\
J1713+0747 & $-15.16$ & $1.6$   &  \cite{NANOGrav12p5yr:2021}\\
J1744-1134 & $-14.16$ & $1.7$   &  \cite{NANOGrav12p5yr:2021}\\
J1747-4036 & $-13.78$ & $3.8$   &  \cite{NANOGrav12p5yr:2021}\\
J1909-3744 & $-14.74$ & $4.05$  &  \cite{GoncharovEtAl:2021}\\
J2145-0750 & $-13.86$ & $1.6$   &  \cite{NANOGrav12p5yr:2021}\\
J2317+1439 & $-16.50$ & $6.0$   &  \cite{NANOGrav12p5yr:2021}\\
\hline
\end{tabular}
\caption{\label{tab:redNoise}
Spin noise parameters for $13$ pulsars. 
[Notes: Values from Ref.~\cite{GoncharovEtAl:2021} are taken directly form their Table 1. 
%Values from Ref.~\cite{NANOGrav12p5yr:2021} are converted from units of ${\rm \mu s \,yr}^{1/2}$ to $\log_{10} \ARed[{\rm yr}^{3/2}]$.]
}
\end{table}
\end{comment}

\textbf{estimating spin noise for other pulsars? }
Spin noise parameters for the remaining pulsars can be estimated using the scaling relation formalism in Shannon and Cordes. 
The spin noise RMS value is given by 
\begin{equation}
\sigma_{\rm SN} = C_2 \mu^{\alpha} | {\dot \mu}_{-15} |^{\beta} t_{\rm yr}^{\gamma}
\end{equation}
where $\mu$......
In this work we use the central fit values for the ${\rm MSP}_{\rm 10,PPTA}$+${\rm NANO}$ ($\{ \ln C_2, \alpha, \beta, \gamma \} = \{-1.2,-0.9,0.8,2.3\}$).

\subsection{JitterNoies}

Some of the pulsars have jitter noise estimates. 
Table~\ref{tab:jitter} shows jitter noise measurements for $13$ MeerTime pulsars from~\cite{ParthasarathyEtAl:2021}.

\begin{table}
\begin{center}
\begin{tabular}{ll}
\hline \hline
Pulsar     & $\sigJitter$ \\
\hline
J0125-2327 & $48  \pm 13$   \\
J0437-4715 & $50  \pm 10$   \\
J0636-3044 & $100 \pm 30$   \\
J0711-6830 & $60  \pm 20$   \\
J1022+1001 & $120 \pm 20$   \\
J1045-4509 & $130 \pm 75$   \\
J1603-7202 & $180 \pm 40$   \\
J1730-2304 & $80  \pm 45$   \\
J1744-1134 & $30  \pm  6$   \\
J1757-5322 & $130 \pm 45$   \\
J1909-3744 & $9   \pm  3$   \\
J2145-0750 & $200 \pm 20$   \\
J2241-5236 & $3.8 \pm  0.8$ \\
\hline\hline
\end{tabular}
\end{center}
\caption{\label{tab:jitter}
Jitter noise measurements.
}
\end{table}
\begin{comment}
\begin{table}
\begin{tabular}{ll}
\hline\hline
Pulsar     & $\sigJitter$ (${\rm ns}$) \\ 
\hline
J0125-2327 & $48  \pm 13$   \\
J0437-4715 & $50  \pm 10$   \\
J0636-3044 & $100 \pm 30$   \\
J0711-6830 & $60  \pm 20$   \\
J1022+1001 & $120 \pm 20$   \\
J1045-4509 & $130 \pm 75$   \\
J1603-7202 & $180 \pm 40$   \\
J1730-2304 & $80  \pm 45$   \\
J1744-1134 & $30  \pm  6$   \\
J1757-5322 & $130 \pm 45$   \\
J1909-3744 & $9   \pm  3$   \\
J2145-0750 & $200 \pm 20$   \\
J2241-5236 & $3.8 \pm  0.8$ \\
\hline\hline
\end{tabular}
\caption{\label{tab:jitter}
Jitter noise measurements from Ref.~\cite{ParthasarathyEtAl:2021}.
}
\end{table}
\end{comment}

\subsection{Current observation strategy}
MeerTime pulsars are observed with a cadence of once every two weeks. 
The total observation time used in each $2$ week observing time is $42592\,{\rm s} = 11\,{\rm hrs}$, which is divided between the MeerTime pulsars. 
The targeted precision for each pulsar is $\sigma = 1 {\rm \mu\, s}$. 

Using the current observation strategy and assuming a GWB signal with $A = 2\times 10^{-15}$ and $\alpha = -2/3$, we can estimate the predicted average SNR for MeerTime from Eqn.~\ref{eqn:snrPIPJPg}. 
After a timespan of XXXX, the average SNR is XXXX
In the following sections we consider whether a re-allocation of the observation time between pulsars could be beneficial to the PTA sensitivity. 


\section{Time Shuffle}
\begin{figure}
\includegraphics[width=0.5\textwidth]{../snr/correlationFunctions.png}
\caption{\label{fig:correlations}
Correlation functions. 
We are using the green function for the shuffle at the moment.
}
\end{figure}

Can we re-organise how the available time is split between pulsars to increase the SNR at$10\,{\rm yr}$?
Shuffle time between pulsar:
\begin{itemize}
\item Choose a pulsar to be the donor pulsar. Subtract $1/4$ of its integration time. 
\item Give time to each other pulsar in turn and compute whether there is an SNR improvement.
\item Repeat for next donor pulsar. 
\end{itemize}

Eqn.~\ref{eqn:hd} has highest values for pulsars that have small angular separations.
But we want good sky coverage to confirm a Hellings-Downs correlation for detection. 
To prevent the time-shuffle algorithm favouring pulsars in the galactic centre, we replace the Hellings and Downs correlation $\chiIJhd$ with a modified angular correlation function $\chi_{IJ}^{\rm mod}$ which corresponds to the difference between a dipole $\chiIJdip$ and the HD correlation
\begin{equation}
\chi_{IJ}^{\rm mod} = \chiIJdip - 2 \chiIJhd
\label{eqn:chimod}
\end{equation}
and $\chiIJdip$ is given by~\cite{AnholmEtAl:2009}
\begin{equation}
\chiIJdip = -\frac{3}{2}
            \left[ \cos \zeta - \frac{4}{3} - 4 \tan^2 \frac{\zeta}{2} 
                   \ln \left( \sin \frac{\zeta}{2} \right) \right]\,.
\end{equation}
The correlation functions are shown in Fig.~\ref{fig:correlations}.




\section{Results}
Two runs using Eqn.~\ref{eqn:chimod} have these results so far: 
\begin{itemize}
\item With red noise (shuffle SNR)/(original SNR): 1.143
\item Without red noise (shuffle SNR)/(original SNR): 1.136
\end{itemize}


\begin{figure}
\includegraphics[width=0.8\textwidth]{../local/integral/quarterTimeShuffleDPHDRed/plots/SNRVTime.png}
\includegraphics[width=0.8\textwidth]{../local/integral/quarterTimeShuffleDPHD/plots/SNRVTime.png}
\caption{\label{fig:SNRVTime}
SNR against PTA lifetime.
The result for the original timings is shown in blue and the shuffle timings in orange.
Top: using red noise parameters in Table~\ref{tab:redNoise}.
Bottom: all red noise parameters set to zero. 
}
\end{figure}

\begin{figure}
\includegraphics[width=0.32\textwidth]{../local/integral/quarterTimeShuffleDPHDRed/plots/obsTimes.png}
\includegraphics[width=0.32\textwidth]{../local/integral/quarterTimeShuffleDPHD/plots/obsTimes.png}
\caption{\label{fig:obsTimes}
Original (blue) and shuffled (orange) observation times for each pulsar.
Left: using red noise parameters in Table~\ref{tab:redNoise}. 
Right: setting all red noise amplitudes to zero.
}
\end{figure}

\begin{figure}
\includegraphics[width=0.9\textwidth]{../local/integral/quarterTimeShuffleDPHDRed/plots/skymapTimeDiff.png}
\includegraphics[width=0.9\textwidth]{../local/integral/quarterTimeShuffleDPHD/plots/skymapTimeDiff.png}
\caption{\label{fig:skymap}
Sky map of shuffle results. 
Orange corresponds to pulsars which gained time and blue to pulsars that lost time. 
The size of the marker is proportional to the amount of time gained/lost. 
The crosses show all pulsar positions. 
Top: with red noise. 
Bottom: zero red noise. 
}
\end{figure}



\section{Conclusions}

The last numbered section should briefly summarise what has been done, and describe
the final conclusions which the authors draw from their work.

\section*{Acknowledgements}

The Acknowledgements section is not numbered. Here you can thank helpful
colleagues, acknowledge funding agencies, telescopes and facilities used etc.
Try to keep it short.
\begin{comment}
Things we have/will use
OzSTAR, hasasia, numpy, scipy, astropy(?), matplotlib
\end{comment}

%%%%%%%%%%%%%%%%%%%%%%%%%%%%%%%%%%%%%%%%%%%%%%%%%%
\section*{Data Availability}

 
The inclusion of a Data Availability Statement is a requirement for articles published in MNRAS. Data Availability Statements provide a standardised format for readers to understand the availability of data underlying the research results described in the article. The statement may refer to original data generated in the course of the study or to third-party data analysed in the article. The statement should describe and provide means of access, where possible, by linking to the data or providing the required accession numbers for the relevant databases or DOIs.




%%%%%%%%%%%%%%%%%%%% REFERENCES %%%%%%%%%%%%%%%%%%

% The best way to enter references is to use BibTeX:

\bibliographystyle{mnras}
\bibliography{ptabib} % if your bibtex file is called example.bib


% Alternatively you could enter them by hand, like this:
% This method is tedious and prone to error if you have lots of references
%\begin{thebibliography}{99}
%\bibitem[\protect\citeauthoryear{Author}{2012}]{Author2012}
%Author A.~N., 2013, Journal of Improbable Astronomy, 1, 1
%\bibitem[\protect\citeauthoryear{Others}{2013}]{Others2013}
%Others S., 2012, Journal of Interesting Stuff, 17, 198
%\end{thebibliography}

%%%%%%%%%%%%%%%%%%%%%%%%%%%%%%%%%%%%%%%%%%%%%%%%%%

%%%%%%%%%%%%%%%%% APPENDICES %%%%%%%%%%%%%%%%%%%%%

\appendix

\section{Some extra material}

If you want to present additional material which would interrupt the flow of the main paper,
it can be placed in an Appendix which appears after the list of references.

%%%%%%%%%%%%%%%%%%%%%%%%%%%%%%%%%%%%%%%%%%%%%%%%%%


% Don't change these lines
\bsp	% typesetting comment
\label{lastpage}
\end{document}

% End of mnras_template.tex
